% !PRINTED 518
\weberpage{518}
3. \gesperrt{An integral functional $\omega$ in the field $\Omega$ is decomposable into prime factors in only one way, if associated prime factors are regarded as not different. Associated functionals contain the same prime factors.}

For let us assume that the two products of prime factors
\weq{6}{\pi_1 \pi_2 \ldots \pi_\nu, \qquad \varkappa_1 \varkappa_2 \ldots \varkappa_\mu}
are associated with one another; then the product $\pi_1 \pi_2 \ldots \pi_\nu$ is divisible by the prime factor $\varkappa_1$, and hence, by \S.\,140, 2., one of the factors, say $\pi_1$, must be divisible by $\varkappa_1$ and consequently be associated with $\varkappa_1$. Then the products
\[
\pi_2 \ldots \pi_\nu, \qquad \varkappa_2 \ldots \varkappa_\mu
\]
are also associated, and consequently one of the factors of the first product, say $\pi_2$, is divisible by $\varkappa_2$ and therefore associated with $\varkappa_2$. In this way one can continue to reason, and it results not only that the number of the $\varkappa$ must agree with the number of the $\pi$, but also that the $\varkappa$ are individually associated with the $\pi$.

Among the prime factors of an integral functional $\omega$ the same one can occur several times, and these mutually equal (or associated) factors can be combined into a power. If $\omega$ is divisible by $\pi^h$, we say that \gesperrt{the prime functional $\pi$ divides $\omega$ $h$ times.}

According to this it has a perfectly definite sense to speak of \gesperrt{the prime factors} of an integral functional. From what has been proved we further conclude:

4. \gesperrt{An integral functional $\alpha$ is divisible by another $\beta$ then and only then, when all prime factors of $\beta$ occur among the prime factors of $\alpha$, and each of them divides $\alpha$ at least as often as it divides $\beta$.}

For if $\alpha$ is divisible by $\beta$, and $\beta$ by $\pi^h$, then $\alpha$ is also divisible by $\pi^h$.

If $\alpha, \beta, \ldots$ are integral functionals, $\pi_1, \pi_2, \ldots$ different prime functionals, and $\varepsilon_1, \varepsilon_2, \ldots$ units, then we can determine exponents $a_1, a_2, \ldots, b_1, b_2, \ldots$ in such a way that
\weq{7}{\begin{aligned}
\varepsilon_1 \alpha &= \pi_1^{a_1} \pi_2^{a_2} \ldots \\
\varepsilon_2 \beta &= \pi_1^{b_1} \pi_2^{b_2} \ldots \\
&\ldots\ldots\ldots\ldots
\end{aligned}}
%CONT>>
